\section{模型评价}

\subsection{误差指标}

本文主要采用 MAE、RMSE 和 WAPE 评价预测误差：
\begin{equation}
MAE=\frac{1}{n}\sum_{i=1}^{n}|y_i-\hat y_i|,
\end{equation}
\begin{equation}
RMSE=\sqrt{\frac{1}{n}\sum_{i=1}^{n}(y_i-\hat y_i)^2},
\end{equation}
\begin{equation}
WAPE=\frac{\sum_{i=1}^{n}|y_i-\hat y_i|}{\sum_{i=1}^{n}|y_i|}.
\end{equation}
MAE 表示平均绝对偏差，便于解释平均误差规模；RMSE 对大误差更敏感，适合检查极端预测偏差；WAPE 用总绝对误差除以总实际销量，适合本题中不同门店和商品销量规模差异较大的情况。对于真实销量为 0 的类别或商品，WAPE 无法定义，需要结合 MAE 和预测合计单独说明。

\subsection{分问题评价}

问题一中，日滚动验证结果显示，门店层面移动平均(7日) WAPE=37.49\%，商品层面简单指数平滑 WAPE=38.14\%，门店--商品层面简单指数平滑 WAPE=74.95\%。补充的 7 日窗口总量验证表明，原问题一主方案在商品层面 WAPE=21.51\%，门店层面 WAPE=22.12\%，门店--商品层面 WAPE=43.06\%。因此，问题一采用简单可解释模型作为基础预测框架。

问题二中，类别聚合后直接预测 + 简单指数平滑的 WAPE=34.74\%，略优于单品预测后按类别加总的 WAPE=34.83\%。销量规模分组 WAPE=24.24\%，相关性聚类 WAPE=31.02\%，但二者不直接回答“同类零食类别预测”要求，因此只作为补充，不作为主方案。

问题三中，合并天气、对数销量和历史控制固定效应回归的验证 WAPE=84.53\%，优于原详细天气固定效应回归的 WAPE=88.09\%。随机森林 grouped weather + history 的 WAPE=84.66\%，可作为预测贡献补充，但不提供清晰方向，也不用于识别因果效应。问题三主结论应写成“条件统计关联”。

问题四中，严格 7 日递推验证是主评价口径。该口径下，问题一门店--商品简单指数平滑 WAPE=76.35\%，综合 Ridge WAPE=76.68\%，14 日滚动均值基准模型 WAPE=77.66\%，7 日移动平均基准模型 WAPE=78.43\%。综合 Ridge 优于问题四内部基准模型，但未超过问题一强基准模型。因此，综合模型的作用主要是融合前三问信息并提供统一解释框架。进一步的低销量鲁棒性检验显示，低销量序列使用指数平滑、常规序列使用综合 Ridge 的混合策略 WAPE=75.85\%，相对问题一强基准降低 0.49 个百分点；该结果说明细粒度稀疏需求需要分层处理，但改进幅度仍应谨慎解释。附加比较还考察了不同低销量阈值、近 28 日均值兜底和层级比例校准；虽然近 28 日均值方案在同一验证集上 WAPE=74.95\%，但留一窗口选模平均改善仅 0.19 个百分点，稳定性证据不足，故不作为正文主预测口径。

\subsection{不确定性}

不确定性主要来自四类来源。第一，未来天气和活动日采用历史同期情景，若实际天气或促销安排偏离该情景，最终预测会变化。第二，低销量门店--商品序列存在大量零销量，当前混合策略下低销量序列整体 WAPE 约为 122.53\%，少量绝对误差即可被低销量分母放大；混合策略虽有所改善，但并未消除稀疏需求带来的不确定性。第三，部分稀有天气样本天数较少，例如雨夹雪、暴雨等，不能把其系数作为稳定主结论。第四，严格递推验证只有 4 个完整 7 日窗口，模型选择结论仍需谨慎。

为辅助解释最终点预测，本文基于当前混合策略在 4 个严格 7 日递推验证窗口上的 7 日聚合残差构造经验区间。对每个门店、商品或类别，先计算“真实 7 日销量减预测 7 日销量”的残差，再取历史残差的 10\% 和 90\% 分位数加到未来点预测上，形成经验 80\% 残差区间。该区间不是严格置信区间，主要反映历史递推误差下的不确定性范围。门店级结果见图 \ref{fig:q4_store_interval}。

\begin{figure}[H]
\centering
\includegraphics[width=0.86\textwidth]{final_forecast_interval_by_store.png}
\caption{基于严格递推残差的门店未来 7 天经验预测区间}
\label{fig:q4_store_interval}
\end{figure}
